\section{Discussion \& Conclusion}
\label{sec:conclusion}

While our ergodic streaming encoder enables bounded, low-latency TTFT, Moonshine v2 still employs a full-attention autoregressive decoder. This means that once the encoder begins emitting features, the decoder must generate tokens one-by-one through a serial loop. For very long transcripts, this sequential generation can add latency to the full output, even though the first tokens appear quickly. Future work could explore monotonic alignment models or streaming-friendly decoding strategies that further reduce end-to-end latency. Additionally, our current models focus exclusively on English ASR. However, the architectural principles of ergodic streaming encoders with sliding-window attention generalize naturally to other languages. Building on our prior work with specialized, language-specific models~\citep{king2025flavors}, we plan to train Moonshine v2 variants for additional languages, enabling low-latency, on-device speech recognition across diverse linguistic contexts.

We introduced Moonshine v2, a family of streaming ASR models designed for latency-critical, on-device applications. By replacing full-attention encoders with ergodic streaming encoders that use sliding-window self-attention, we achieve bounded TTFT independent of utterance length while maintaining strong transcription accuracy. Our models achieve state-of-the-art results on standard benchmarks, matching the performance of models 6x their size while running significantly faster. These results demonstrate that carefully designed local attention can rival the accuracy of global attention at a fraction of the computational cost, making real-time, interactive speech interfaces practical on resource-constrained edge devices.
